\documentclass[10pt,a4paper]{amsart}
\usepackage[margin=19mm]{geometry}
\usepackage{amsmath,amssymb,mathtools}
\usepackage[scr=rsfs,cal=euler]{mathalfa}
\usepackage{xfrac}
\usepackage[unicode,hidelinks,bookmarks=false]{hyperref}
\newcommand{\ie}{i.e.\ }
\newcommand{\cf}{cf.\ }
\newcommand{\resp}{respectively\ }
\newcommand{\Subsection}{\subsection}
\providecommand{\nobreakdash}{\nobreak\hbox{-}}
\setlength{\emergencystretch}{2em}
\multlinegap0pt
\usepackage{amssymb}
\usepackage{float}
\usepackage{dsfont}
\usepackage{stackengine}
\usepackage{enumerate}
\usepackage{xcolor}
\usepackage[normalem]{ulem}
\usepackage{cancel}
\newtheorem{lemma}{Lemma}
\newcommand{\bq}{\begin{equation}}
\newcommand{\calN}{\mathcal N}
\newcommand{\eq}{\end{equation}}
\newcommand{\intr}{\int_{\mR^3}}
\newcommand{\lal}{\langle}
\newcommand{\lt}{\left}
\newcommand{\mR}{\mathbb{R}}
\newcommand{\mfu}{\mathfrak{u}}
\newcommand{\ral}{\rangle}
\newcommand{\rdp}{\textnormal{d}p}
\newcommand{\rt}{\right}
\newcommand{\vrh}{\varrho}
\newcommand{\vth}{\vartheta}
\title{Incompressible Navier--Stokes--Fourier limit from a nonlinear quantum Fokker--Planck equation}
\author{Young-Pil Choi, Byung-Hoon Hwang, Ju-Hwan Hyun}
\date{Source: arXiv:2607.27583v1 (CC BY 4.0)}
\begin{document}
\maketitle

\noindent\emph{Source and licence.} Incompressible Navier--Stokes--Fourier limit from a nonlinear quantum Fokker--Planck equation. Version arXiv:2607.27583v1 (CC BY 4.0), distributed by arXiv under the Creative Commons Attribution 4.0 International licence, http://creativecommons.org/licenses/by/4.0/, as recorded in the arXiv licence metadata for this exact version and shown on the abstract page https://arxiv.org/abs/2607.27583v1. Full licence notice: \texttt{licenses/arxiv-2607.27583-CC-BY-4.0.txt}.\par\smallskip

\noindent\textit{Editorial scope.} Counted unit: the article's lemma lem:tr-flux (source0.tex lines 1730--1743). The article's complete original proof of it is delivered with it (source0.tex lines 1745--1785, one \texttt{proof} environment of 41 lines). No mathematics is written, repaired or re-derived: the statement and the proof are the article's own lines, in the article's own notation, and no retained line is substituted or re-flowed.  Retained uncounted as the source-local context the proof needs: lines 1261--1293; lines 1333--1344.  Nothing was omitted from any retained span, so the package carries no omission record.  No retained line contains a citation, so the wrapper carries no bibliography and no bibliography entry.  No retained line contains a figure environment, an image-inclusion command or drawing code, so no image is delivered and the package declares no asset dependency.  Editorial formatting only, in addition: an AMS \texttt{amsart} preamble that restates the article's own declarations and macros used by the retained lines, namely the environments \texttt{lemma} on separate counters, together with the standard packages those lines need.  The retained environments print in this document as Lemma 1 in order of appearance, whereas the article prints the same items with the numbering of its own full document.  The complete native source of the article as distributed in the arXiv e-print for this exact version is bundled unchanged as \texttt{sources/206356/source0.tex}.
\begin{lemma}\label{lem:AB-solv}
There exist unique functions $\widetilde A \in \mathbb{H}^{3\times3}$, $\widetilde B \in \mathbb{H}^3$ such that
\bq\label{eq:AB-aux}
L\widetilde A = A, \quad L\widetilde B = B
\eq
in the weak sense. More precisely, after extending $\mathfrak{B}$ continuously to $\mathbb{H}$, the identities in \eqref{eq:AB-aux} mean that
\bq\label{eq:AB-weak}
\mathfrak{B}(\widetilde A,\Phi)=-\lal A,\Phi\ral_{L^2_p}
\quad\text{for all } \Phi\in\mathbb{H}^{3\times3},
\quad
\mathfrak{B}(\widetilde B,\Psi)=-\lal B,\Psi\ral_{L^2_p}
\quad\text{for all } \Psi\in\mathbb{H}^{3}.
\eq

Moreover, there exists a constant $C_\hbar>0$ such that
\bq\label{eq:AB-D}
| \widetilde A |_D + | \widetilde B |_D \le C_\hbar,
\eq
and
\bq \label{eq:AB-L2}
\nabla_p\widetilde A, \; \nabla_p\widetilde B, \; p\widetilde A, \; p\widetilde B \in L^2_p.
\eq

There exist radial scalar functions $\mathfrak a=\mathfrak a(|p|)$ and $\mathfrak b=\mathfrak b(|p|)$
such that
\[
\widetilde A(p) = \mathfrak a(|p|) \lt( p\otimes p - \frac{|p|^2}{3} \mathbb I_3 \rt),\quad \widetilde B(p) = \mathfrak b(|p|) p.
\]
Finally,
\[
\lal \widetilde A, A \ral_{L^2_p} < 0, \quad \lal \widetilde B, B \ral_{L^2_p} < 0.
\]
\end{lemma}

We define the quantum viscosity coefficient by
\bq \label{eq:nu-h}
\nu_\hbar:=-\frac{1}{10m_2 }\lal\widetilde A ,A \ral_{L^2_p}>0.
\eq
We also introduce
\[
\chi_\hbar:=-\frac{2}{3m_0\beta\phi }\lal\widetilde B ,B \ral_{L^2_p}>0
\]
and define the quantum thermal diffusivity by
\[
\kappa_\hbar:=\frac{2\chi_\hbar}{(1+\beta )\phi }=-\frac{4}{3(1+\beta )m_0\beta\phi ^2}\lal\widetilde B ,B \ral_{L^2_p}>0.
\]

\begin{lemma}\label{lem:tr-flux}
Let $g=g(x,p)$ and write
\[
P g=\lt\{\frac{\vrh}{m_0 }+\frac{\mfu\cdot p}{m_2 }+\frac{\vth}{m_0\beta}\lt(\frac{|p|^2}{\phi }-1\rt)\rt\}\sqrt{\mu_\hbar}.
\]
Then,
\bq\label{eq:A-tr-app}
\lal p\cdot\nabla_xP g,\widetilde A \ral_{L^2_p}=-\nu_\hbar\Sigma(\mfu)
\eq
and
\bq\label{eq:B-tr-app}
\lal p\cdot\nabla_xP g,\widetilde B \ral_{L^2_p}=-\chi_\hbar\nabla_x\vth.
\eq
\end{lemma}

\begin{proof}
We first consider the stress tensor. By the representation of $\widetilde A$ in Lemma \ref{lem:AB-solv}, $\widetilde A$ is even in $p$. Thus, the density and thermal components of $Pg$ do not contribute to $\lal p\cdot\nabla_xPg,\widetilde A\ral_{L^2_p}$ since the corresponding integrands are odd functions of $p$. Hence only the momentum component contributes, and
\bq\label{eq:A-tr-comp}
\lal p\cdot\nabla_xP g,\widetilde A \ral_{L^2_p}=\frac{1}{m_2 }\sum_{i,j=1}^3\partial_{x_i}\mfu_j\intr p_ip_j\widetilde A \sqrt{\mu_\hbar}\,\rdp.
\eq
Since $L $ is rotationally invariant and $A $ is a symmetric traceless tensor, the same properties hold for $\widetilde A $. Therefore,
\bq\label{eq:A-rot}
\intr p_ip_j(\widetilde A )_{k\ell}\sqrt{\mu_\hbar}\,\rdp=\frac{1}{10}\lal\widetilde A ,A \ral_{L^2_p}\lt(\delta_{ik}\delta_{j\ell}+\delta_{i\ell}\delta_{jk}-\frac{2}{3}\delta_{ij}\delta_{k\ell}\rt).
\eq
Substituting \eqref{eq:A-rot} into \eqref{eq:A-tr-comp} and using the definition \eqref{eq:nu-h}, we obtain
\[
\lal p\cdot\nabla_xP g,\widetilde A \ral_{L^2_p}=-\nu_\hbar\lt\{\nabla_x\mfu+(\nabla_x\mfu)^T-\frac{2}{3}(\nabla_x\cdot\mfu) \mathbb I_3\rt\},
\]
which proves \eqref{eq:A-tr-app}.

We next consider the heat flux. By the representation of $\widetilde B$ in Lemma \ref{lem:AB-solv}, $\widetilde B$ is odd in $p$. Thus, the momentum component of $Pg$ does not contribute since the corresponding velocity integral is proportional to
\[
\intr p_i p_k(\widetilde B)_j\sqrt{\mu_\hbar}\,\rdp,
\]
whose integrand is odd in $p$. The density component is not excluded by this symmetry; instead, it vanishes by the orthogonality $\widetilde B\in\calN^\perp$, namely
\[
\intr p_i(\widetilde B )_j\sqrt{\mu_\hbar}\,\rdp=0.
\]
Hence, only the thermal component contributes:
\[
\lal p\cdot\nabla_xP g,\widetilde B \ral_{L^2_p}=\frac{1}{m_0\beta}\sum_{i,j=1}^3\partial_{x_i}\vth\intr p_i\lt(\frac{|p|^2}{\phi }-1\rt)(\widetilde B )_j\sqrt{\mu_\hbar}\,\rdp.
\]
Using again the orthogonality of $\widetilde B $ to $p\sqrt{\mu_\hbar}$, we obtain
\[
\intr p_i\lt(\frac{|p|^2}{\phi }-1\rt)(\widetilde B )_j\sqrt{\mu_\hbar}\,\rdp=\frac{1}{\phi }\intr p_i|p|^2(\widetilde B )_j\sqrt{\mu_\hbar}\,\rdp.
\]
By rotational symmetry,
\[
\intr p_i|p|^2(\widetilde B )_j\sqrt{\mu_\hbar}\,\rdp=\frac{2}{3}\lal\widetilde B ,B \ral_{L^2_p}\delta_{ij}.
\]
It follows that
\[
\lal p\cdot\nabla_xP g,\widetilde B \ral_{L^2_p}=\frac{2}{3m_0\beta\phi }\lal\widetilde B ,B \ral_{L^2_p}\nabla_x\vth=-\chi_\hbar\nabla_x\vth,
\]
which proves \eqref{eq:B-tr-app}.
\end{proof}

\end{document}
